% GENERATED FILE. Edit canonical lessons, then run npm run generate:books.
% canonical-source-hash: fnv1a64:44d4a674e1eba4df
% canonical-lessons: RU-C324-pobeda, RU-C324-pobezhdat, RU-C324-voyevat, RU-C324-pogibat, RU-C324-osvobozhdat

\chapter{Victory}
\label{ch:ru-victory}
\hlchaptermodality{324}

\begin{chapteropening}
\textbf{By the end of this chapter:} \emph{I can talk about a victory, winning, waging war, dying in battle and setting people free.}

Complete the last lesson of chapter 324: I can talk about a victory, winning, waging war, dying in battle and setting people free.
\end{chapteropening}

\section[\texorpdfstring{\ru{победа} (pobéda)}{победа (pobéda)}]{\ru{победа} (pobéda) — a victory, a win}

\label{lesson:RU-C324-pobeda}



\begin{quote}
Before the new one: say the Russian for to fight, then the Russian for an enemy.
\end{quote}

\subsection*{You'll want to know: \ru{победа}}
\textbf{\ru{победа}} (\emph{pobéda}) — \textquotedblleft{}a victory, a win\textquotedblright{}.

\begin{quote}
\ru{Победа} \ru{была} \ru{очень} \ru{важной} \ru{для} \ru{всех}.
\end{quote}

\begin{quote}
\emph{Pobéda bylá óchen vázhnoy dlya vsekh.}
\end{quote}

\begin{quote}
The victory was very important for everyone.
\end{quote}

Feminine. \emph{\ru{День} \ru{Победы}} — Victory Day, 9 May.

\begin{cousinweb}
\emph{\ru{победа}} is inherited from Old East Slavic \emph{poběda} \textquotedblleft{}victory, defeat\textquotedblright{} (EtymDB), built on \emph{\ru{по}-} and \emph{\ru{беда}} \textquotedblleft{}misfortune\textquotedblright{}: first the trouble one side brings on the other.
\end{cousinweb}

\begin{grammarlens}[title={What you've built}]
One more word for telling, explaining and describing at B1.
\end{grammarlens}

\subsection*{Your turn}
\begin{itemize}
\raggedright
  \item \emph{Say it:} \emph{pobéda}
  \item \emph{Say it:} \emph{pobéda}, once more
  \item \emph{Say it:} a sentence of your own with \emph{pobéda}
  \item \emph{Recall:} say the Russian for to fight, then the Russian for an enemy, then say \emph{pobéda} again
\end{itemize}

\begin{tcolorbox}[breakable,colback=teal!4,colframe=teal!35!black,title={Before you move on}]
What does \ru{победа} mean? (\textquotedblleft{}A victory, a win\textquotedblright{}.) Say it once more.
\end{tcolorbox}
\section[\texorpdfstring{\ru{побеждать} (pobezhdát)}{побеждать (pobezhdát)}]{\ru{побеждать} (pobezhdát) — to win, to defeat}

\label{lesson:RU-C324-pobezhdat}



\begin{quote}
Before the new one: say the Russian for an enemy, then the Russian for a victory.
\end{quote}

\subsection*{You'll want to know: \ru{побеждать}}
\textbf{\ru{побеждать}} (\emph{pobezhdát}) — \textquotedblleft{}to win, to defeat\textquotedblright{}.

\begin{quote}
\ru{Наша} \ru{команда} \ru{победила} \ru{в} \ru{финале}.
\end{quote}

\begin{quote}
\emph{Násha kománda pobedíla v finále.}
\end{quote}

\begin{quote}
Our team won the final.
\end{quote}

Imperfective; perfective \emph{\ru{победить}} (\emph{pobedít}), which has no \emph{I} form in the future: Russians say \emph{\ru{я} \ru{одержу} \ru{победу}}.

\begin{cousinweb}
\emph{\ru{побеждать}} is built on \emph{\ru{победа}} \textquotedblleft{}victory\textquotedblright{}, Old East Slavic \emph{poběda} (EtymDB), in the Church Slavonic shape \emph{-\ru{жд}-}.
\end{cousinweb}

\begin{grammarlens}[title={What you've built}]
One more word for telling, explaining and describing at B1.
\end{grammarlens}

\subsection*{Your turn}
\begin{itemize}
\raggedright
  \item \emph{Say it:} \emph{pobezhdát}
  \item \emph{Say it:} \emph{pobezhdát}, once more
  \item \emph{Say it:} a sentence of your own with \emph{pobezhdát}
  \item \emph{Recall:} say the Russian for an enemy, then the Russian for a victory, then say \emph{pobezhdát} again
\end{itemize}

\begin{tcolorbox}[breakable,colback=teal!4,colframe=teal!35!black,title={Before you move on}]
What does \ru{побеждать} mean? (\textquotedblleft{}To win, to defeat\textquotedblright{}.) Say it once more.
\end{tcolorbox}
\section[\texorpdfstring{\ru{воевать} (voyevát)}{воевать (voyevát)}]{\ru{воевать} (voyevát) — to wage war, to fight}

\label{lesson:RU-C324-voyevat}



\begin{quote}
Before the new one: say the Russian for a victory, then the Russian for to win.
\end{quote}

\subsection*{You'll want to know: \ru{воевать}}
\textbf{\ru{воевать}} (\emph{voyevát}) — \textquotedblleft{}to wage war, to fight\textquotedblright{}.

\begin{quote}
\ru{Его} \ru{дедушка} \ru{воевал} \ru{четыре} \ru{года}.
\end{quote}

\begin{quote}
\emph{Yevó dédushka voyevál chetýre góda.}
\end{quote}

\begin{quote}
His grandfather fought in the war for four years.
\end{quote}

Imperfective; present \emph{\ru{воюю}}, \emph{\ru{воюешь}}. \emph{\ru{Воевать} \ru{с}} + instrumental.

\begin{cousinweb}
\emph{\ru{воевать}} is built on \emph{\ru{воин}} \textquotedblleft{}a warrior\textquotedblright{}, from Proto-Slavic \emph{vojinŭ}, with the verb ending \emph{-\ru{евать}} (EtymDB).
\end{cousinweb}

\begin{grammarlens}[title={What you've built}]
One more word for telling, explaining and describing at B1.
\end{grammarlens}

\subsection*{Your turn}
\begin{itemize}
\raggedright
  \item \emph{Say it:} \emph{voyevát}
  \item \emph{Say it:} \emph{voyevát}, once more
  \item \emph{Say it:} a sentence of your own with \emph{voyevát}
  \item \emph{Recall:} say the Russian for a victory, then the Russian for to win, then say \emph{voyevát} again
\end{itemize}

\begin{tcolorbox}[breakable,colback=teal!4,colframe=teal!35!black,title={Before you move on}]
What does \ru{воевать} mean? (\textquotedblleft{}To wage war, to fight\textquotedblright{}.) Say it once more.
\end{tcolorbox}
\section[\texorpdfstring{\ru{погибать} (pogibát)}{погибать (pogibát)}]{\ru{погибать} (pogibát) — to die (in war or an accident), to perish}

\label{lesson:RU-C324-pogibat}



\begin{quote}
Before the new one: say the Russian for to win, then the Russian for to wage war.
\end{quote}

\subsection*{You'll want to know: \ru{погибать}}
\textbf{\ru{погибать}} (\emph{pogibát}) — \textquotedblleft{}to die (in war or an accident), to perish\textquotedblright{}.

\begin{quote}
\ru{В} \ru{той} \ru{войне} \ru{погибло} \ru{много} \ru{людей}.
\end{quote}

\begin{quote}
\emph{V toy voyné pogíblo mnógo lyudéy.}
\end{quote}

\begin{quote}
Many people died in that war.
\end{quote}

Imperfective; perfective \emph{\ru{погибнуть}} (\emph{pogíbnut}), past \emph{\ru{погиб}}. For a violent death, not an illness.

\begin{cousinweb}
\emph{\ru{погибать}} is built on \emph{\ru{по}-} and \emph{\ru{гибнуть}} \textquotedblleft{}to perish\textquotedblright{}, inherited from Proto-Slavic \emph{gybnǫti} (EtymDB), the root of \emph{\ru{гибель}} \textquotedblleft{}doom\textquotedblright{}.
\end{cousinweb}

\begin{grammarlens}[title={What you've built}]
One more word for telling, explaining and describing at B1.
\end{grammarlens}

\subsection*{Your turn}
\begin{itemize}
\raggedright
  \item \emph{Say it:} \emph{pogibát}
  \item \emph{Say it:} \emph{pogibát}, once more
  \item \emph{Say it:} a sentence of your own with \emph{pogibát}
  \item \emph{Recall:} say the Russian for to win, then the Russian for to wage war, then say \emph{pogibát} again
\end{itemize}

\begin{tcolorbox}[breakable,colback=teal!4,colframe=teal!35!black,title={Before you move on}]
What does \ru{погибать} mean? (\textquotedblleft{}To die (in war or an accident), to perish\textquotedblright{}.) Say it once more.
\end{tcolorbox}
\section[\texorpdfstring{\ru{освобождать} (osvobozhdát)}{освобождать (osvobozhdát)}]{\ru{освобождать} (osvobozhdát) — to free, to liberate; to vacate}

\label{lesson:RU-C324-osvobozhdat}



\begin{quote}
Before the new one: say the Russian for to wage war, then the Russian for to die.
\end{quote}

\subsection*{You'll want to know: \ru{освобождать}}
\textbf{\ru{освобождать}} (\emph{osvobozhdát}) — \textquotedblleft{}to free, to liberate; to vacate\textquotedblright{}.

\begin{quote}
\ru{Солдаты} \ru{освободили} \ru{город}.
\end{quote}

\begin{quote}
\emph{Soldáty osvobodíli górod.}
\end{quote}

\begin{quote}
The soldiers liberated the town.
\end{quote}

Imperfective; perfective \emph{\ru{освободить}} (\emph{osvobodít}), \emph{\ru{освобожу}}. Also to free up: \emph{\ru{освободить} \ru{место}}.

\begin{cousinweb}
\emph{\ru{освобождать}} is built on \emph{\ru{о}-} and \emph{\ru{свобода}} \textquotedblleft{}freedom\textquotedblright{}, in the Church Slavonic shape \emph{-\ru{жд}-}; EtymDB lists \emph{\ru{освободить}} with no parent.
\end{cousinweb}

\begin{grammarlens}[title={What you've built}]
One more word for telling, explaining and describing at B1.
\end{grammarlens}

\subsection*{Your turn}
\begin{itemize}
\raggedright
  \item \emph{Say it:} \emph{osvobozhdát}
  \item \emph{Say it:} \emph{osvobozhdát}, once more
  \item \emph{Say it:} a sentence of your own with \emph{osvobozhdát}
  \item \emph{Recall:} say the Russian for to wage war, then the Russian for to die, then say \emph{osvobozhdát} again
\end{itemize}

\begin{tcolorbox}[breakable,colback=teal!4,colframe=teal!35!black,title={Before you move on}]
What does \ru{освобождать} mean? (\textquotedblleft{}To free, to liberate; to vacate\textquotedblright{}.) Say it once more.
\end{tcolorbox}
